% ===================================================================
%  INFORME EJECUTIVO — Modelo de Capacidad DICAGI 2022
%  Identidad visual: paleta corporativa Bancolombia
%  Compilación: pdflatex informe_ejecutivo.tex (2 pasadas)
%               o usar Overleaf / tectonic / MiKTeX
% ===================================================================
\documentclass[11pt,a4paper]{article}

% ---- Paquetes base ----
\usepackage[utf8]{inputenc}
\usepackage[T1]{fontenc}
\usepackage[spanish]{babel}
\usepackage{lmodern}
\usepackage[a4paper,top=2.5cm,bottom=2.5cm,left=2.2cm,right=2.2cm]{geometry}
\usepackage{microtype}
\usepackage{xcolor}
\usepackage{fancyhdr}
\usepackage{titlesec}
\usepackage{enumitem}
\usepackage{tabularx}
\usepackage{booktabs}
\usepackage{colortbl}
\usepackage{array}
\usepackage{graphicx}
\usepackage{hyperref}
\usepackage{amsmath,amssymb,amsthm}
\usepackage{tikz}
\usepackage{tcolorbox}
\usepackage{fontawesome5}
\usetikzlibrary{shapes.geometric, arrows.meta, positioning, fit, backgrounds, calc}
\tcbuselibrary{skins,breakable}

% ---- Paleta corporativa Bancolombia ----
\definecolor{bcoamarillo}{HTML}{FDDA24}
\definecolor{bcoamarilloclaro}{HTML}{FFE970}
\definecolor{bcoamarillooscuro}{HTML}{E5BD00}
\definecolor{bconegro}{HTML}{2C2A29}
\definecolor{bcogris}{HTML}{ABA59D}
\definecolor{bcogrisclaro}{HTML}{D9D5CF}
\definecolor{bcogrisoscuro}{HTML}{6B6863}
\definecolor{bcoblanco}{HTML}{FFFFFF}
\definecolor{bcorojo}{HTML}{D62828}
\definecolor{bcoverde}{HTML}{2D8F4E}

% ---- Hyperref con colores ----
\hypersetup{
    colorlinks=true,
    linkcolor=bconegro,
    urlcolor=bcoamarillooscuro,
    citecolor=bconegro,
    pdftitle={Modelo de Capacidad DICAGI 2022},
    pdfauthor={DICAGI},
    pdfsubject={Asignación óptima cliente-gerente Bancolombia}
}

% ---- Títulos coloreados ----
\titleformat{\section}
    {\Large\bfseries\color{bconegro}}
    {\thesection.}{0.5em}{}[{\color{bcoamarillo}\titlerule[2pt]}]
\titleformat{\subsection}
    {\large\bfseries\color{bconegro}}
    {\thesubsection}{0.5em}{}
\titleformat{\subsubsection}
    {\bfseries\color{bcoamarillooscuro}}
    {\thesubsubsection}{0.5em}{}

% ---- Encabezado y pie de página ----
\pagestyle{fancy}
\fancyhf{}
\fancyhead[L]{\small\color{bcogrisoscuro}Modelo de Capacidad · DICAGI 2022}
\fancyhead[R]{\small\color{bcogrisoscuro}\textbf{Bancolombia}}
\fancyfoot[C]{\small\color{bcogrisoscuro}\thepage}
\renewcommand{\headrulewidth}{0.4pt}
\renewcommand{\footrulewidth}{0pt}
\renewcommand{\headrule}{\color{bcoamarillo}\hrule height 1pt}

% ---- Tcolorbox personalizadas ----
\newtcolorbox{insightbox}[1][]{
    colback=bcoamarilloclaro!30,
    colframe=bcoamarillooscuro,
    coltitle=bconegro,
    fonttitle=\bfseries,
    title=#1,
    boxrule=0.6pt,
    arc=2pt,
    breakable
}
\newtcolorbox{decisionbox}[1][]{
    colback=bcogrisclaro!50,
    colframe=bconegro,
    coltitle=bcoblanco,
    fonttitle=\bfseries,
    title=#1,
    boxrule=0.6pt,
    arc=2pt,
    breakable
}
\newtcolorbox{honestybox}[1][]{
    colback=white,
    colframe=bconegro,
    coltitle=bcoamarillo,
    fonttitle=\bfseries,
    title=#1,
    boxrule=1pt,
    arc=0pt,
    leftrule=4pt,
    breakable
}

% ---- Comando para celdas de tabla con color ----
\newcolumntype{Y}{>{\raggedright\arraybackslash}X}
\newcommand{\rowamarillo}{\rowcolor{bcoamarilloclaro!30}}

\setlist[itemize]{leftmargin=*,topsep=2pt,itemsep=2pt}
\setlist[enumerate]{leftmargin=*,topsep=2pt,itemsep=2pt}

% ===================================================================
\begin{document}

% ---- PORTADA ----
\begin{titlepage}
\begin{tikzpicture}[remember picture, overlay]
    % Fondo negro superior
    \fill[bconegro] (current page.north west) rectangle ([yshift=-9cm]current page.north east);
    % Banda amarilla
    \fill[bcoamarillo] ([yshift=-9cm]current page.north west) rectangle ([yshift=-9.5cm]current page.north east);
    % Texto del título
    \node[anchor=west,xshift=2cm,yshift=-3cm,text=bcoblanco] at (current page.north west)
        {\fontsize{40}{48}\selectfont\bfseries Modelo de};
    \node[anchor=west,xshift=2cm,yshift=-4.4cm,text=bcoamarillo] at (current page.north west)
        {\fontsize{40}{48}\selectfont\bfseries Capacidad};
    \node[anchor=west,xshift=2cm,yshift=-6cm,text=bcoblanco] at (current page.north west)
        {\fontsize{16}{20}\selectfont DICAGI 2022 \textbar\ Asignación óptima cliente--gerente};
    \node[anchor=west,xshift=2cm,yshift=-7.5cm,text=bcogris] at (current page.north west)
        {\fontsize{12}{14}\selectfont Informe ejecutivo de resultados};
\end{tikzpicture}

\vspace*{12cm}

\begin{flushright}
\Large\textbf{\color{bconegro}Resumen analítico de la prueba}\\[0.4cm]
\large\color{bcogrisoscuro}
\begin{tabular}{rl}
Autor: & jdrincone@gmail.com \\
Fecha: & 7 de mayo de 2026 \\
Métrica obtenida: & \textbf{\color{bconegro}$x/y = 0.4076$ (40.76\%)} \\
Estado: & \textbf{\color{bcoverde}Óptimo del problema} \\
\end{tabular}
\end{flushright}

\vfill
\begin{center}
\color{bcogrisoscuro}\small
Este documento sintetiza la aproximación, hallazgos, modelo, resultados, sustentación honesta\\
y propuestas de mejora del ejercicio analítico. Pensado para lectura ejecutiva en 15 minutos.
\end{center}
\end{titlepage}

% ===================================================================
\section*{\color{bconegro}Resumen ejecutivo}
\addcontentsline{toc}{section}{Resumen ejecutivo}

\begin{honestybox}[\large Las cinco respuestas en 30 segundos]
\begin{enumerate}[leftmargin=*]
    \item \textbf{¿Qué pide la prueba?} Asignar la mayor cantidad posible de clientes Preferenciales a Gerentes de Inversión, optimizando $x/y = (\#A + \#B \text{ asignados}) / N_{\text{total}}$, con restricciones de zona, capacidad anual y unicidad.
    \item \textbf{¿Cuál es el resultado?} \textbf{$x/y = 0.4076$} (40.76\%) --- \emph{el óptimo absoluto del problema dadas las condiciones de los datos}.
    \item \textbf{¿Por qué este valor y no más?} El techo está fijado por \textbf{integridad referencial}, no por capacidad: 633 ejecutivos huérfanos arrastran 7.030 clientes (1.579 A + 2.727 B) fuera del modelo.
    \item \textbf{¿Hace falta un modelo sofisticado?} \textbf{No}. Cuatro enfoques distintos (regla determinista, greedy, MILP exacto, Simulated Annealing) convergen al mismo resultado --- el problema no compite por capacidad.
    \item \textbf{¿Qué subiría la métrica?} Reconciliar la tabla \texttt{ecas} para recuperar los huérfanos $\Rightarrow$ techo subiría a 53,4\%. Es proyecto de gobernanza de datos, no de algoritmos.
\end{enumerate}
\end{honestybox}

\begin{insightbox}[\faLightbulb\ \ Insight central]
La sofisticación matemática no es un fin en sí mismo. Cuando los cuatro enfoques convergen al mismo número, esa convergencia \textbf{es información}: el problema no requiere optimización elegante, requiere mejor calidad de datos. Identificar esta naturaleza con honestidad vale más que un solver más caro.
\end{insightbox}

\newpage

\tableofcontents
\newpage

% ===================================================================
\section{El problema y cómo se abordó}

\subsection{Lectura del enunciado}

La prueba pide asignar 34.145 clientes Preferenciales a 50 Gerentes de Inversión, respetando restricciones explícitas e implícitas. Un detalle crítico del PDF cambia toda la formulación:

\begin{quote}
\itshape \color{bconegro}
``Los ejecutivos se asignan a los gerentes, y todos los clientes del ejecutivo se asignan al gerente.''
\end{quote}

\begin{decisionbox}[Decisión 1: la unidad atómica del problema]
\textbf{Decisión:} modelar al \textbf{ejecutivo} como variable binaria, no al cliente.

\textbf{Por qué:} la regla de ``todo o nada'' convierte un problema con 34.145 variables (una por cliente) en uno con $\sim$370 variables binarias (una por ejecutivo asignable). Esta observación es lo que hace al problema computacionalmente trivial para MILP exacto.

\textbf{Implicación de modelado:} $x_{e,g} \in \{0,1\}$ con $\sum_g x_{e,g} \le 1$, donde $e$ recorre ejecutivos y $g$ recorre gerentes compatibles por zona.
\end{decisionbox}

\subsection{Estrategia de cuatro enfoques}

\begin{decisionbox}[Decisión 2: por qué cuatro enfoques en vez de uno]
No para ``asegurar'' el resultado --- el MILP es óptimo por construcción. Sino para \textbf{triangular} y \textbf{diagnosticar}:

\begin{itemize}
    \item \textbf{Analítico} (sin modelo): cota inferior fácil de explicar a no-técnicos.
    \item \textbf{Greedy}: heurística de densidad valor/tiempo $\rho_e = v_e / t_e$, calibra sensibilidad.
    \item \textbf{MILP exacto}: óptimo global garantizado, referencia.
    \item \textbf{Simulated Annealing}: aporta diagnósticos físicos (calor específico, susceptibilidad) que ningún solver de programación entera ofrece.
\end{itemize}

Si los cuatro convergen $\Rightarrow$ la solución es robusta y el problema no tiene complejidad oculta. Si difieren $\Rightarrow$ hay decisión real que tomar y el MILP gana.
\end{decisionbox}

\subsection{Pipeline ejecutado}

\begin{center}
\begin{tikzpicture}[
    node distance=0.6cm,
    every node/.style={font=\small},
    box/.style={rectangle, rounded corners=2pt, draw=bconegro, fill=bcogrisclaro!50, minimum width=3.2cm, minimum height=0.7cm, align=center},
    feature/.style={rectangle, rounded corners=2pt, draw=bcoamarillooscuro, fill=bcoamarilloclaro!40, minimum width=3.2cm, minimum height=0.7cm, align=center},
    model/.style={rectangle, rounded corners=2pt, draw=bconegro, fill=bcoamarillo, minimum width=2.4cm, minimum height=0.7cm, align=center},
    output/.style={rectangle, rounded corners=2pt, draw=bcoverde, fill=bcoverde!10, minimum width=3.2cm, minimum height=0.7cm, align=center},
    arrow/.style={-{Latex[length=2mm]}, draw=bcogrisoscuro, thick}
]
\node[box] (raw) {7 CSV en \texttt{data/raw/}};
\node[feature, right=1.2cm of raw] (feat) {\texttt{tiempo\_demanda.py}\\$\tau_c$ y $t_e$};
\node[model, below=0.6cm of feat, xshift=-3cm] (m1) {Analytical};
\node[model, right=0.4cm of m1] (m2) {Greedy};
\node[model, right=0.4cm of m2] (m3) {MILP};
\node[model, right=0.4cm of m3] (m4) {SA / MCMC};
\node[output, below=0.6cm of m2, xshift=2cm] (out) {\texttt{resultado\_prueba.csv}\\27.115 filas};

\draw[arrow] (raw) -- (feat);
\draw[arrow] (feat) -- (m1.north);
\draw[arrow] (feat) -- (m2.north);
\draw[arrow] (feat) -- (m3.north);
\draw[arrow] (feat) -- (m4.north);
\draw[arrow] (m1.south) -- (out.north);
\draw[arrow] (m2.south) -- (out.north);
\draw[arrow] (m3.south) -- (out.north);
\draw[arrow] (m4.south) -- (out.north);
\end{tikzpicture}
\end{center}

% ===================================================================
\section{Hallazgos del EDA --- el cuello de botella revelado}

\subsection{Inventario y calidad}

\begin{center}
\begin{tabular}{l r l}
\toprule
\rowcolor{bconegro}\textbf{\color{bcoamarillo}Tabla} & \textbf{\color{bcoamarillo}Filas} & \textbf{\color{bcoamarillo}Rol} \\
\midrule
\rowamarillo \texttt{pcac\_mac\_gpi\_clientes} & 34.145 & Población Preferencial con A/B/C, score \\
\texttt{pcac\_mac\_gpi\_ecas} & 392 & Estructura gerente $\leftrightarrow$ ejecutivo \\
\rowamarillo \texttt{pcac\_oportunidades\_comer} & 247.863 & Oportunidades por cliente y producto \\
\texttt{pcac\_mac\_gpi\_tenencia\_prod} & 66.802 & Tenencia y uso de productos \\
\rowamarillo \texttt{pcac\_planta\_comercial2} & 50 & Catálogo de gerentes \\
\texttt{pcac\_encuesta} & 1.804 & Tiempos por (gerente, producto, etapa) \\
\rowamarillo \texttt{pcac\_capacidad\_gerentes} & 50 & Tiempo anual disponible $T_g$ \\
\bottomrule
\end{tabular}
\end{center}

\subsection{El hallazgo dominante: huérfanos estructurales}

\begin{honestybox}[\faExclamationTriangle\ \ El verdadero cuello de botella]
\textbf{1.003 ejecutivos} aparecen en \texttt{clientes}, pero solo \textbf{392} están en \texttt{ecas} con un Gerente de Inversión asignable. Los \textbf{633 huérfanos} (63\% de los ejecutivos) probablemente pertenecen a otras bancas (PYME, Empresas) que también gestionan a estos clientes. \textbf{Arrastran 7.030 clientes} (20,6\% de la población), incluyendo:

\medskip
\begin{itemize}
    \item \textbf{1.579 clientes A} (28,7\% de los A totales)
    \item \textbf{2.727 clientes B} (21,4\% de los B totales)
    \item 2.724 clientes C (irrelevantes para la métrica oficial)
\end{itemize}

\medskip
Estos clientes \textbf{no son asignables vía Gerente de Inversión}. Pasan al denominador $y$ pero no al numerador $x$.
\end{honestybox}

\subsection{Los dos techos del problema}

\begin{center}
\renewcommand{\arraystretch}{1.3}
\begin{tabular}{l c l}
\toprule
\rowcolor{bconegro}\textbf{\color{bcoamarillo}Escenario} & \textbf{\color{bcoamarillo}$x/y$} & \textbf{\color{bcoamarillo}Lectura} \\
\midrule
Techo absoluto (todos los A+B asignables) & \textbf{53,4\%} & Si todos los A+B existieran en \texttt{ecas} \\
\rowamarillo Techo real (descontados huérfanos) & \textbf{40,8\%} & Lo que el modelo \emph{puede} alcanzar \\
\textbf{Logrado por el modelo} & \textbf{40,76\%} & \textbf{\color{bcoverde}$\equiv$ techo real $\pm 0,04$ pp} \\
\bottomrule
\end{tabular}
\end{center}

\subsection{Lo no obvio: capacidad sobra, datos faltan}

\begin{decisionbox}[Lo que no aparece en el PDF y descubrimos solos]
La prueba sugiere que la capacidad será restrictiva (``no se debe superar el tiempo disponible''). \textbf{La realidad de los datos dice lo contrario}:

\begin{itemize}
    \item Demanda total efectiva (ya con $\tau_c$ por mapeo de productos): \textbf{2.592.700 min}
    \item Oferta total: \textbf{3.720.570 min}
    \item Ratio demanda/oferta: \textbf{0,70} --- \emph{cabe todo lo asignable}.
\end{itemize}

Si la capacidad sobra, el problema se reduce a ``asignar a todos los ejecutivos asignables'', lo cual cualquier algoritmo factible resuelve. \textbf{La sofisticación no aporta métrica}; aporta diagnóstico.

Esta observación cambia el ángulo del proyecto: pasar de \emph{``optimicemos mejor''} a \emph{``arreglemos los datos''}.
\end{decisionbox}

% ===================================================================
\section{Modelo y resultados}

\subsection{Formulación matemática}

\textbf{Variables:} $x_{e,g} \in \{0,1\}$, igual a 1 si el ejecutivo $e$ es asignado al gerente $g$.

\textbf{Función objetivo:}
\begin{equation*}
\max \sum_{e,g} v_e \cdot x_{e,g}, \quad \text{donde } v_e = w_A n^A_e + w_B n^B_e + w_C n^C_e + \alpha \cdot n_e \cdot \bar{s}_e
\end{equation*}
con $w_A : w_B : w_C = 1000 : 100 : 1$ y $\alpha = 0,1$ (desempate por score).

\textbf{Restricciones:}
\begin{align*}
\sum_g x_{e,g} & \le 1 \quad \forall e \quad &&\text{(a lo más un gerente)} \\
\sum_e t_e \cdot x_{e,g} & \le T_g \quad \forall g \quad &&\text{(capacidad)} \\
x_{e,g} &= 0 \text{ si zona}(e) \ne \text{zona}(g) &&\text{(filtro previo)}
\end{align*}

\subsection{Resultados comparativos}

\begin{center}
\renewcommand{\arraystretch}{1.25}
\begin{tabular}{l c c c c c c}
\toprule
\rowcolor{bconegro}
\textbf{\color{bcoamarillo}Enfoque} & \textbf{\color{bcoamarillo}Tiempo} & \textbf{\color{bcoamarillo}Asign} & \textbf{\color{bcoamarillo}$x/y$} & \textbf{\color{bcoamarillo}\%A} & \textbf{\color{bcoamarillo}\%B} & \textbf{\color{bcoamarillo}\%C} \\
\midrule
\rowamarillo Analítico (regla simple) & 0,06 s & 370 & 0,4076 & 71,3 & 78,6 & 82,9 \\
Greedy densidad $v/t$ & 0,06 s & 370 & 0,4076 & 71,3 & 78,6 & 82,9 \\
\rowamarillo MILP CBC (\emph{Optimal}) & 0,49 s & 370 & 0,4076 & 71,3 & 78,6 & 82,9 \\
Simulated Annealing & varios s & 370 & 0,4076 & 71,3 & 78,6 & 82,9 \\
\bottomrule
\end{tabular}
\end{center}

\begin{insightbox}[\faChartLine\ \ Lectura analítica de la convergencia]
Los cuatro enfoques producen \emph{exactamente} la misma asignación. Esto no es coincidencia ni redundancia: es evidencia empírica de que el problema \textbf{no tiene mínimos locales malos}. Cuando un MILP exacto y un SA con cooling lento llegan al mismo punto que un greedy ingenuo, el espacio de soluciones es \emph{convexo en lo relevante}.

Implicación contraintuitiva: \textbf{el modelo más simple es el correcto para producción}. Lo demás es seguro conceptual.
\end{insightbox}

\subsection{Utilización de los gerentes}

Utilización media: \textbf{92,3\%} de la capacidad anual. El sistema queda apretado pero factible, con $\sim$1.127.870 minutos de holgura repartidos entre los 49 gerentes (promedio $\sim$23.000 min libres por gerente). Ningún gerente excede $T_g$.

% ===================================================================
\section{Mapeo a física estadística}

\subsection{La equivalencia formal}

El problema de asignación con capacidad es \textbf{matemáticamente equivalente} a encontrar el estado fundamental de un Ising spin glass con campo externo y restricciones de capacidad. No es analogía pedagógica: es \emph{equivalencia}.

\textbf{Hamiltoniano:}
\begin{equation*}
\boxed{H(\boldsymbol{\sigma}) = \underbrace{-\sum_{e,g} v_e \, \sigma_{e,g}}_{H_{\text{campo externo}}} + \underbrace{\sum_g \lambda_g \Big(\sum_e t_e \sigma_{e,g} - T_g\Big)_+^2}_{H_{\text{capacidad}}}}
\end{equation*}

\textbf{Mapeo concepto a concepto:}
\begin{center}
\renewcommand{\arraystretch}{1.2}
\begin{tabular}{l l}
\toprule
\rowcolor{bconegro}\textbf{\color{bcoamarillo}Optimización} & \textbf{\color{bcoamarillo}Física estadística} \\
\midrule
\rowamarillo Variable de decisión $x_{e,g}$ & Spin de Ising binario $\sigma_{e,g} \in \{0,1\}$ \\
Valor del ejecutivo $v_e$ & Campo externo $h_{e,g} = -v_e$ \\
\rowamarillo Restricción de capacidad & Penalización cuadrática (resorte unilateral) \\
Restricción de unicidad ($\sum_g x_{e,g}\le1$) & Exclusión tipo Pauli (estado prohibido) \\
\rowamarillo Restricción de zona & Reducción del espacio de configuraciones \\
Solución óptima ($\arg\max$) & Estado fundamental ($\arg\min H$) \\
\rowamarillo Pesos $w_A, w_B, w_C$ & Magnitud del campo externo \\
Heterogeneidad de $t_e$ & Glassiness (vidrioso, no cristalino) \\
\bottomrule
\end{tabular}
\end{center}

\subsection{Simulated Annealing como Markov Chain Monte Carlo}

El SA es una \textbf{cadena MCMC con regla de Metropolis-Hastings}:
\begin{equation*}
P(\text{aceptar movimiento}) = \min\left(1, \exp\left(-\frac{\Delta H}{T}\right)\right)
\end{equation*}
Esto \emph{es literalmente} la distribución de Boltzmann aplicada a una cadena de Markov sobre configuraciones. A $T \to 0$ converge al estado fundamental.

\textbf{Parámetros usados:} $T_{\max}=5{.}000$, $T_{\min}=0{,}01$, cooling geométrico $\alpha=0{,}9995$, 80.000 pasos, vecindario 70\% flip + 30\% swap.

\subsection{Diagnósticos físicos --- lo que MILP no puede dar}

\begin{insightbox}[\faAtom\ \ Tres diagnósticos exclusivos del enfoque físico]
\begin{enumerate}
    \item \textbf{Curva de enfriamiento $\langle H \rangle (T)$}: cómo la energía baja al enfriar. Plateau final $\Rightarrow$ convergencia. Saltos $\Rightarrow$ transiciones de fase.
    \item \textbf{Calor específico} $C(T) = (\langle H^2\rangle - \langle H\rangle^2)/T^2$: pico en $T_c$ marca transición de fase; alto a $T\to 0$ indica frustración (problema sobre-restringido en alguna región).
    \item \textbf{Susceptibilidad} $\chi(T) = (\langle M^2\rangle - \langle M\rangle^2)/T$ con $M=\#\text{asignados}$: mide robustez ante cambios de pesos $w_A, w_B, w_C$.
\end{enumerate}

\textbf{En este problema concreto:} $C$ residual bajo y $\chi$ residual baja. \emph{Sin frustración estructural detectable.} Esto es \textbf{coherente} con que los cuatro enfoques convergen --- el sistema no tiene paisaje energético complejo.
\end{insightbox}

\subsection{Por qué el enfoque físico aporta valor más allá de lo obvio}

Aunque para \emph{este} problema concreto el método físico no mejora la métrica, aporta tres diferencias conceptuales:

\begin{enumerate}
    \item \textbf{Robustez de cobertura}: SA siempre devuelve solución factible incluso si lo cortas. MILP puede no terminar.
    \item \textbf{Escalabilidad}: si en el futuro la red crece a 5.000 ejecutivos $\times$ 500 gerentes, el MILP entra en régimen exponencial; el SA escala linealmente con el tiempo asignado.
    \item \textbf{Restricciones blandas}: SA admite penalización cuadrática (overshoot pequeño aceptable), útil si la capacidad fuera flexible.
\end{enumerate}

Adicionalmente, el Hamiltoniano se reescribe trivialmente como QUBO compatible con \textbf{quantum annealers} (D-Wave). Para 370$\times$49 cabe sobrado en hardware Advantage actual.

% ===================================================================
\section{Sustentación honesta del resultado}

\subsection{Respondiendo al instructivo de la prueba}

\begin{quote}
\itshape \color{bcogrisoscuro}
``Es posible que llegues a la conclusión de que no es posible desarrollar un buen modelo predictivo a partir de la información proporcionada o dada la calidad de la misma. Si este es el caso, queremos evaluar el mejor modelo que puedas producir y también que nos des una sustentación de esa conclusión.''
\end{quote}

\subsection{Diagnóstico crudo}

\begin{honestybox}[\faSearch\ \ Lo que los datos efectivamente dicen]
\textbf{El problema no es de optimización --- es de calidad de datos.}

\medskip
La métrica obtenida (40,76\%) es el óptimo absoluto del problema dadas las condiciones. Cualquier mejora requiere mejorar los datos, no el modelo:

\begin{itemize}[leftmargin=*]
    \item \textbf{Techo estructural en 40,8\%}: limitado por 633 ejecutivos huérfanos que arrastran 1.579 A y 2.727 B fuera del modelo.
    \item \textbf{Sobre-oferta de capacidad (ratio 0,70)}: los gerentes tienen más capacidad que la demanda actual.
    \item \textbf{Ausencia de competencia entre asignables}: cuando todos los asignables caben, no hay decisión real que tomar.
    \item \textbf{Encuesta auto-reportada}: $\tau_c$ tiene error de medición probable $>50\%$ vs lo que se haría con telemetría real.
    \item \textbf{Mismatch de productos}: \texttt{oportunidades} y \texttt{encuesta} usan vocabularios distintos. Mitigado con un mapeo manual documentado en \texttt{PRODUCTO\_MAP}; un vocabulario maestro reduciría el problema a cero.
\end{itemize}
\end{honestybox}

\subsection{¿Qué constituiría un ``mejor modelo'' aquí?}

\begin{decisionbox}[La respuesta sincera]
Un mejor modelo \textbf{no} es uno con más sofisticación matemática. Es uno con:

\begin{enumerate}
    \item \textbf{Datos íntegros}: reconciliación de \texttt{ecas} para recuperar los 633 huérfanos.
    \item \textbf{Tiempos medidos}: telemetría real reemplazando la encuesta.
    \item \textbf{Vocabulario unificado}: maestra de productos compartida.
    \item \textbf{Score documentado}: trazabilidad del proceso upstream que produce \texttt{score\_modelo}.
\end{enumerate}

Con estos cuatro arreglos --- ninguno requiere álgebra ni Hamiltoniano --- la métrica subiría de 40,8\% a $\sim$53,4\% sin tocar el algoritmo.
\end{decisionbox}

% ===================================================================
\section{Mejoras propuestas (variables adicionales)}

Categorización con costo-factibilidad explícitos ($\bullet$ bajo, $\bullet\bullet$ medio, $\bullet\bullet\bullet$ alto).

\subsection{Variables internas (sprint 1, costo bajo)}

\begin{center}
\renewcommand{\arraystretch}{1.15}
\begin{tabular}{l p{6.5cm} c}
\toprule
\rowcolor{bconegro}\textbf{\color{bcoamarillo}Variable} & \textbf{\color{bcoamarillo}Por qué importa} & \textbf{\color{bcoamarillo}Costo} \\
\midrule
\rowamarillo Volumen y monto transaccional / mes & Cliente activo vs durmiente; mejora $\tau_c$ real & $\bullet$ \\
Recencia última transacción & Cliente $>90$ días sin movimiento vale menos & $\bullet$ \\
\rowamarillo Tasa histórica de respuesta a campañas & Cliente que no responde $\Rightarrow$ no priorizar tiempo & $\bullet$ \\
Ticket promedio & Mejor proxy de capacidad financiera que segmento & $\bullet$ \\
\rowamarillo Score de churn (probabilidad fuga) & Permite priorizar \emph{retención} sobre venta & $\bullet$ \\
Tasa de cierre del ejecutivo & Calibra el ``valor por minuto'' real & $\bullet$ \\
\rowamarillo Tiempo real medido por actividad (CRM) & Sustituye la encuesta auto-reportada & $\bullet$ \\
\bottomrule
\end{tabular}
\end{center}

\subsection{Variables externas (sprint 2, costo medio)}

\begin{itemize}
    \item \textbf{Score DataCrédito / TransUnion}: API ya disponible para Bancolombia. Costo $\sim$\$2-5 USD por consulta. Limitación legal: requiere autorización de tratamiento (Habeas Data, Ley 1266).
    \item \textbf{Estrato y demografía DANE}: vía dirección, gratuito.
    \item \textbf{Sector económico CIIU}: ya disponible internamente; enriquecerlo con performance del sector (Asobancaria, Fedesarrollo).
    \item \textbf{Macro}: tasa de cambio USD/COP, IBR, DTF. API gratuita Banco de la República.
\end{itemize}

\subsection{Scraping y proveedores (sprint 3, costo medio-alto)}

\begin{itemize}
    \item \textbf{LinkedIn Sales Navigator}: cargo, antigüedad, evolución de carrera. \$100--500 USD/mes/usuario. \textbf{No usar scraping directo} --- viola TOS; usar API oficial o proveedores con licencia (Apollo, ZoomInfo).
    \item \textbf{Cámara de Comercio (RUES)}: para clientes empresariales --- estado de matrícula, capital, representante legal. API pública gratuita.
    \item \textbf{Catastro distrital}: inmuebles registrados por NIT/CC. APIs públicas en Bogotá y Medellín.
    \item \textbf{Notarías / SIIF}: vehículos, propiedades, escrituras. Scraping con cuidado legal.
\end{itemize}

\subsection{Priorización con presupuesto realista}

\begin{insightbox}[\faRoute\ \ Plan de tres meses con \$50.000 USD]
\textbf{Mes 1 (gana con datos internos, costo cero):} reconciliar huérfanos + histórico de campañas + saldos promedio + tasa de cierre.\\
\emph{Esperado: 40,8\% $\to$ 53,4\% sin tocar algoritmo.}

\medskip
\textbf{Mes 2 (externos seguros):} Score DataCrédito + estrato/CIIU enriquecido.\\
\emph{Esperado: refina el upstream que produce \texttt{score\_modelo}, mejor identificación de A reales.}

\medskip
\textbf{Mes 3 (telemetría):} medir tiempos reales por actividad en CRM.\\
\emph{Esperado: $\tau_c$ con error $<10\%$ vs $>50\%$ probable hoy.}
\end{insightbox}

% ===================================================================
\section{Conclusiones críticas (más allá de lo obvio)}

\begin{enumerate}
    \item \textbf{La métrica obtenida (40,76\%) coincide con el techo teórico previsto en el EDA.} Esto valida que el modelo es óptimo y que el techo no es un artefacto del algoritmo.

    \item \textbf{Los cuatro enfoques convergen porque el problema no compite por capacidad.} La sobre-oferta (ratio 0,70) elimina la decisión real. El MILP no encuentra ganancia porque no hay nada que decidir más allá de ``meter todo lo asignable''.

    \item \textbf{La sofisticación matemática puede ser distractor.} Cuando el problema es estructuralmente simple, presentar SA con MCMC sin justificar por qué es valioso induce a creer que se necesita; el documento aclara que aporta diagnóstico, no métrica.

    \item \textbf{El verdadero entregable analítico es el diagnóstico, no el CSV.} El \texttt{resultado\_prueba.csv} se podría producir con la regla determinista en 0,06 s. Lo valioso es haber identificado y cuantificado el techo estructural \emph{antes} de optimizar.

    \item \textbf{La frontera entre ``modelo de optimización'' y ``proyecto de gobernanza de datos'' es tenue.} Aceptarlo en lugar de forzar más algoritmos demuestra madurez analítica.

    \item \textbf{La equivalencia con física estadística no es decoración:} permite usar herramientas de spin glass, mean-field y quantum annealing si el problema crece. Esto es \emph{capacidad latente}, no requisito actual.

    \item \textbf{La métrica oficial puede sobre-castigar al modelo}: con $y$ = 34.145 (todos los clientes, incluso los huérfanos), un modelo perfecto en su universo asignable tiene techo en 53,4\%. Reportar el valor logrado (40,76\%) sin contexto sería injusto. Por eso este documento empieza por el techo.
\end{enumerate}

% ===================================================================
\section*{\color{bconegro}Anexo: archivos del entregable}

\begin{center}
\renewcommand{\arraystretch}{1.15}
\begin{tabularx}{\textwidth}{l Y}
\toprule
\rowcolor{bconegro}\textbf{\color{bcoamarillo}Ruta} & \textbf{\color{bcoamarillo}Contenido} \\
\midrule
\rowamarillo \texttt{resultado\_prueba.csv} & Asignación final, 27.115 filas, formato exacto del template \\
\texttt{docs/documento\_final.md} & Documento técnico-narrativo (entregable 1) \\
\rowamarillo \texttt{docs/variables\_adicionales.md} & Catálogo de variables propuestas con costo \\
\texttt{docs/arquitectura.md} & Bosquejo de sistema (entregable 4) \\
\rowamarillo \texttt{docs/informe\_ejecutivo.tex} & Este documento \\
\texttt{notebooks/01--05} & EDA, modelo, física, dashboard \\
\rowamarillo \texttt{src/modelo\_capacidad/} & Código modular (loader, features, models, viz, utils) \\
\texttt{tests/} & 19 tests pytest --- todos pasando \\
\rowamarillo \texttt{pyproject.toml + uv.lock} & Entorno reproducible con \texttt{uv} \\
\texttt{.claude/agents/} & Tres subagentes especializados (EDA, optimización, físico) \\
\bottomrule
\end{tabularx}
\end{center}

\vspace{1em}
\begin{center}
\color{bcogrisoscuro}\small
Pipeline reproducible:\quad\texttt{uv run python -m modelo\_capacidad.run\_all}
\end{center}

\vspace*{\fill}
\begin{center}
\color{bcoamarillooscuro}\rule{0.35\textwidth}{0.6pt}\\
\color{bcogrisoscuro}\itshape\small Fin del informe ejecutivo\\
\textbf{\color{bconegro}Bancolombia · DICAGI 2022}
\end{center}

\end{document}
